\documentclass[12pt]{article}
\usepackage[utf8]{inputenc}
\usepackage[T1]{fontenc}
\usepackage{hyperref}
\title{Sakachep Language Resources for K-12 Curriculum Development}
\author{Compiled by Claude Sonnet 5.5 (claude-sonnet-5-5) for Andr\'as Kornai}
\date{}
\begin{document}
\maketitle

\section{Introduction}
Sakachep (also Khelma, Achep; ISO 639-3 sch; Glottocode saka1283) is a Central Kuki-Chin-Mizo language of Sino-Tibetan stock spoken in northeast India, mainly in Assam, Manipur and Nagaland, with small communities reported in Mizoram and Tripura. Speaker estimates conflict: about 25,000 (Wikipedia, 2003), about 21,000 in the Joshua Project ethnic total, and 10,000 in the Global Recordings Network listing. Dialect names given include Khelma, Achep, Mar-Achep, Thang Achep and Sak Achep. In searches I found no evidence of schooling in the language and no standardized orthography, dictionary or grammar. The material below is therefore thin.

\section{Primary Educational Resources}
No primers, textbooks or school programmes for Sakachep turned up in searches. Nothing here should be read as an existing school resource.

\section{Linguistic and Reference Materials}
\begin{itemize}
\item Wikipedia, ``Sakachep language'': a general overview with classification and dialect names. \url{https://en.wikipedia.org/wiki/Sakachep_language}
\item Wikipedia, ``Khelma people'': background on the related Khelma community, described as predominantly Presbyterian. \url{https://en.wikipedia.org/wiki/Khelma_people}
\item Numerals comparison page (Max Planck, EVA) listing Sakachep, Assam and Manipur: a small linguistic data item. \url{https://lingweb.eva.mpg.de/channumerals/Sakachep.htm}
\end{itemize}
No dictionary or descriptive grammar of Sakachep was found.

\section{Religious and Community Resources}
\begin{itemize}
\item Scripture Earth page for Sakachep. It lists a New Testament (``Pathien Lekhabu Inthieng Chongtung Thar'', YouVersion code SAKNT), Gospel of Mark and Luke films, and Global Recordings audio. Its Bible text is the largest body of written Sakachep and could serve as a source of orthographic practice. \url{https://www.scriptureearth.org/00eng.php?iso=sch}
\item Joshua Project profile of the Sakachep in India (population and religion data). \url{https://www.joshuaproject.net/index.php/people_groups/21559}
\end{itemize}

\section{Digital and Audio Resources}
\begin{itemize}
\item Global Recordings Network, Sakachep (sch): audio recordings. \url{https://www.globalrecordings.net/en/language/sch}
\end{itemize}

\section{Curriculum Development Context}
Sakachep is reported as stable and used by all in the community, but it is not known to be taught in schools. A curriculum would have to begin with community work on orthography (the Bible translation uses a Latin-based writing), a word list and basic literacy primers. Existing resources are limited to scripture, a few audio items and general encyclopedic descriptions. Speaker numbers and the presence of the language in Tripura are poorly documented and should be verified directly with the community.
\end{document}